\section{Full Theoretical Proofs}
\label{app:full_theoretical_proofs}

This appendix proves the results in Section~\ref{sec:theoretical_analysis} under explicit conditions. Self-supervised calibration does not guarantee score separation, a general GRU need not be contractive, and encoder-input normalization does not bound raw aggregation updates. For analysis, $\mathcal{A}_t=\mathcal{G}_t\mathbin{\dot\cup}\mathcal{B}_t$ partitions protocol-valid clients. Monitoring does not remove otherwise valid clients; empty-cohort rounds leave the model unchanged and are excluded from the effective-round index.

%======================================================================
\subsection{Benign Trajectory Stability}
\label{app:proof_benign_stability}

\noindent\textit{Proof of Proposition~\ref{prop:benign_representation_stability}.} Fix a benign client and suppress its client index where no confusion arises. The participation index $n$ advances only when that client submits a protocol-valid update. Write $\mathbf{h}_n=\Phi_{\phi_n}(\mathbf{h}_{n-1},\mathbf{x}_n)$ and $\mathbf{z}_n=\Psi_{\phi_n}(\mathbf{h}_n)$. The historical hidden state is retained from the preceding participation and is detached for differentiation, but detachment does not change its numerical value. Let $d_n=\|\mathbf{h}_n-\mathbf{h}_{n-1}\|_2$ for $n\geq1$. Since $\mathbf{h}_{n-1}=\Phi_{\phi_{n-1}}(\mathbf{h}_{n-2},\mathbf{x}_{n-1})$, the first inequality of Eq.~\eqref{eq:theory_encoder_regularities} yields, for $n\geq2$,
\begin{align}
d_n&=\|\Phi_{\phi_n}(\mathbf{h}_{n-1},\mathbf{x}_n)-\Phi_{\phi_{n-1}}(\mathbf{h}_{n-2},\mathbf{x}_{n-1})\|_2\nonumber\\
&\leq\chi d_{n-1}+L_x\|\mathbf{x}_n-\mathbf{x}_{n-1}\|_2+L_\phi\|\phi_n-\phi_{n-1}\|_2\nonumber\\
&\leq\chi d_{n-1}+C_h,
\label{eq:app_hidden_recursion}
\end{align}
where $C_h:=L_xD_{\mathrm{ben}}+L_\phi D_\phi$. Here $D_\phi$ bounds the full encoder change between consecutive observed participations of this client, including updates made during its absence. Repeated substitution gives
\begin{equation}
d_n\leq\chi^{n-1}d_1+C_h\sum_{j=0}^{n-2}\chi^j=\chi^{n-1}d_1+C_h\frac{1-\chi^{n-1}}{1-\chi}.
\label{eq:app_hidden_finite_bound}
\end{equation}
Since $0\leq\chi<1$, the first term vanishes, so
\begin{equation}
\limsup_{n\to\infty}d_n\leq\frac{C_h}{1-\chi}.
\label{eq:app_hidden_limit}
\end{equation}
The projection-head regularity in Eq.~\eqref{eq:theory_encoder_regularities} then implies
\begin{equation}
\begin{aligned}
\|\mathbf{z}_n-\mathbf{z}_{n-1}\|_2&\leq L_zd_n+L_\psi\|\phi_n-\phi_{n-1}\|_2\\
&\leq L_zd_n+L_\psi D_\phi.
\end{aligned}
\label{eq:app_latent_difference}
\end{equation}
Taking the limit superior and applying Eq.~\eqref{eq:app_hidden_limit} proves
\begin{equation}
\limsup_{n\to\infty}\|\mathbf{z}_n-\mathbf{z}_{n-1}\|_2\leq\frac{L_zC_h}{1-\chi}+L_\psi D_\phi=B_z,
\end{equation}
as required by Eq.~\eqref{eq:latent_stability_main}.

For the angular statement, suppose $\|\mathbf{z}_n\|_2\geq m_z\geq\varepsilon_{\cos}>0$ for all sufficiently large $n$. The stabilized cosine in Section~\ref{subsec:method_multi_resolution} uses $\max(\|\mathbf{z}\|_2,\varepsilon_{\cos})$ in each denominator. Hence, under the stated norm condition, it equals ordinary cosine similarity. For any nonzero $\mathbf{a},\mathbf{b}$, the exact identity
\begin{equation}
\|\mathbf{a}-\mathbf{b}\|_2^2=(\|\mathbf{a}\|_2-\|\mathbf{b}\|_2)^2+2\|\mathbf{a}\|_2\|\mathbf{b}\|_2[1-\cos(\mathbf{a},\mathbf{b})]
\label{eq:app_cosine_identity}
\end{equation}
gives $1-\cos(\mathbf{a},\mathbf{b})\leq\|\mathbf{a}-\mathbf{b}\|_2^2/(2m_z^2)$ whenever both norms exceed $m_z$. Substituting successive latent states, taking the limit superior, and using Eq.~\eqref{eq:latent_stability_main} proves Eq.~\eqref{eq:latent_cosine_stability}. \hfill$\square$

\noindent The contraction assumption, the between-participation encoder-drift bound, and the positive latent-norm bound are sufficient conditions rather than automatic consequences of recurrent architecture or reconstruction training. With frozen encoder parameters, $D_\phi=0$, and Eq.~\eqref{eq:latent_stability_main} specializes to $B_z=L_zL_xD_{\mathrm{ben}}/(1-\chi)$. The limit-superior statement concerns successive valid participations, not successive global rounds.

%======================================================================
\subsection{Fast--Slow Reference Separation}
\label{app:proof_fast_slow}

\noindent\textit{Proof of Lemma~\ref{lem:fast_slow_separation}.} Consider $0\leq n\leq N$, with $s^n=s^0+\nu_{\mathrm{drift}}n$. For $\gamma\in(0,1)$, the scalar EMA is $\bar s_{\gamma}^n=\gamma s^n+(1-\gamma)\bar s_{\gamma}^{n-1}$, initialized at $\bar s_{\gamma}^0=s^0$. Define $e_{\gamma}^n=s^n-\bar s_{\gamma}^n$. Substituting the EMA recursion and $s^n-s^{n-1}=\nu_{\mathrm{drift}}$ gives
\begin{equation}
e_{\gamma}^n=(1-\gamma)\big(\nu_{\mathrm{drift}}+e_{\gamma}^{n-1}\big),\qquad e_{\gamma}^0=0.
\label{eq:app_ema_lag_recursion}
\end{equation}
Induction on $n$ yields the exact finite-step expression
\begin{equation}
e_{\gamma}^n=\nu_{\mathrm{drift}}\sum_{k=1}^{n}(1-\gamma)^k.
\label{eq:app_ema_finite_lag}
\end{equation}
Since $\bar s_{\gamma}^n=s^n-e_{\gamma}^n$, we have
\begin{align}
\bar s_{\mathrm{fast}}^n-\bar s_{\mathrm{slow}}^n&=e_{\mathrm{slow}}^n-e_{\mathrm{fast}}^n\nonumber\\
&=\nu_{\mathrm{drift}}\sum_{k=1}^{n}\big[(1-\gamma_{\mathrm{slow}})^k-(1-\gamma_{\mathrm{fast}})^k\big].
\end{align}
Every summand is strictly positive because $0<1-\gamma_{\mathrm{fast}}<1-\gamma_{\mathrm{slow}}<1$ and $\nu_{\mathrm{drift}}>0$. This establishes Eq.~\eqref{eq:fast_slow_main}. \hfill$\square$

\noindent The scalar projection of a vector EMA obeys the same recursion by linearity. Nevertheless, the statement concerns separation of projected fast and slow references, not guaranteed nonzero $\Delta_{\mathrm{MR}}$: angular distances can coincide even for different references. The formula holds on a finite drift interval; it does not assume that a uniformly bounded representation follows unbounded linear drift indefinitely. In the implemented update order, current-round features are evaluated using the anchors from the previous valid participation.

%======================================================================
\subsection{Bounded Harmful Aggregation Influence}
\label{app:proof_harmful_influence}

\noindent\textit{Proof of Theorem~\ref{thm:bounded_harmful_influence}.} Let $h_i^t=q_i^t\exp(\widetilde r_i^t)>0$. Since $q_i^t=\sigma(\widetilde r_i^t)$, we have $\exp(\widetilde r_i^t)=q_i^t/(1-q_i^t)$ and $h_i^t=s(q_i^t)$, where $s(q)=q^2/(1-q)$. Moreover, $s'(q)=q(2-q)/(1-q)^2>0$ on $(0,1)$.

Define $H_t=\sum_{c_i\in\mathcal{B}_t}h_i^t$ and $Q_t=\sum_{c_i\in\mathcal{G}_t}h_i^t$. Score separation gives
\begin{equation}
H_t\leq |\mathcal{B}_t|s(q_{\mathrm{bad}})=:\overline H_t,\qquad Q_t\geq |\mathcal{G}_t|s(q_{\mathrm{ben}})=:\underline Q_t.
\label{eq:app_harmful_unnormalized_bounds}
\end{equation}
For nonempty cohorts, $H_t+Q_t>0$. The normalized weights of Eq.~\eqref{eq:interaction_aware_weight} satisfy
\begin{equation}
\sum_{c_i\in\mathcal{B}_t}\alpha_i^t=\frac{H_t}{H_t+Q_t}\leq\frac{\overline H_t}{\overline H_t+\underline Q_t}.
\label{eq:app_harmful_mass}
\end{equation}
Here $x/(x+y)$ increases with $x$ and decreases with $y$ for nonnegative masses. The bound also covers all-benign and all-harmful cohorts. By the triangle inequality and the assumed raw-update limit,
\begin{equation}
\|E_{\mathrm{bad}}^t\|_2\leq\sum_{c_i\in\mathcal{B}_t}\alpha_i^t\|\Delta w_i^t\|_2\leq G_{\mathrm{bad}}\frac{\overline H_t}{\overline H_t+\underline Q_t}=A_t,
\end{equation}
proving Eq.~\eqref{eq:harmful_bound_main}. \hfill$\square$

\noindent The proof assumes, but does not establish, benign--harmful score separation. With an all-harmful cohort, $A_t=G_{\mathrm{bad}}$: reliability weighting cannot reduce its total mass. Bounded logits do not bound $\|\Delta w_i^t\|_2$; finite parameter-space influence requires an independent raw-update limit, enforced only if the implementation includes one.

%======================================================================
\subsection{Convergence to a Reliability-Controlled Neighborhood}
\label{app:proof_stationarity}

\noindent\textit{Proof of Theorem~\ref{thm:trajrep_stationarity}.} Index nonempty-cohort rounds by $t=0,\ldots,T-1$. Set $g_t=\nabla F(w^t)$, $G_t=G_{\mathrm{ben}}^t+E_{\mathrm{bad}}^t$, and $w^{t+1}=w^t+\eta G_t$. By $L$-smoothness,
\begin{equation}
F(w^{t+1})\leq F(w^t)+\eta\langle g_t,G_t\rangle+\frac{L\eta^2}{2}\|G_t\|_2^2.
\label{eq:app_smoothness_step}
\end{equation}
Let $\mathbb E_t[\cdot]=\mathbb E[\cdot\mid\mathcal F_t]$. Using $\|u+v\|_2^2\leq2\|u\|_2^2+2\|v\|_2^2$, the measurability of $g_t,F(w^t)$, and Eq.~\eqref{eq:benign_aggregate_conditions} yields
\begin{equation}
\begin{aligned}
\mathbb E_t[F(w^{t+1})]\leq{}&F(w^t)-\eta a\|g_t\|_2^2\\
&+L\eta^2(b_1\|g_t\|_2^2+b_2\sigma^2)\\
&+\eta\mathbb E_t[\langle g_t,E_{\mathrm{bad}}^t\rangle]\\
&+L\eta^2\mathbb E_t[\|E_{\mathrm{bad}}^t\|_2^2].
\end{aligned}
\label{eq:app_conditional_descent}
\end{equation}
Young's inequality gives $\langle g_t,E_{\mathrm{bad}}^t\rangle\leq(a/2)\|g_t\|_2^2+(1/(2a))\|E_{\mathrm{bad}}^t\|_2^2$. Together with $\|E_{\mathrm{bad}}^t\|_2\leq A_t$, this implies
\begin{equation}
\begin{aligned}
\mathbb E_t[F(w^{t+1})]\leq{}&F(w^t)\\
&-\eta\left(\frac a2-L\eta b_1\right)\|g_t\|_2^2\\
&+L\eta^2b_2\sigma^2\\
&+\eta\left(\frac{1}{2a}+L\eta\right)\mathbb E_t[A_t^2].
\end{aligned}
\label{eq:app_descent_recursion}
\end{equation}
Since $0<\eta\leq a/(4Lb_1)$, $a/2-L\eta b_1\geq a/4$. Taking total expectations and telescoping yields
\begin{equation}
\begin{aligned}
\frac{\eta a}{4}\sum_{t=0}^{T-1}\mathbb E\|\nabla F(w^t)\|_2^2\leq{}&F(w^0)-\mathbb E[F(w^T)]\\
&+TL\eta^2b_2\sigma^2\\
&+\eta\left(\frac{1}{2a}+L\eta\right)\sum_{t=0}^{T-1}\mathbb E[A_t^2].
\end{aligned}
\label{eq:app_telescoped_bound}
\end{equation}
Using $F(w^T)\geq F_{\inf}$ and dividing by $T\eta a/4$ yields Eq.~\eqref{eq:stationarity_main}, with $\overline{A^2}_T=T^{-1}\sum_t\mathbb E[A_t^2]$. \hfill$\square$

\noindent The result concerns average stationarity for stationary $F$ under expected-descent and second-moment conditions on the weighted benign aggregate. Those conditions may fail when weighting suppresses useful directions, benign participation is sparse, or adversaries dominate. Neither encoder training nor this proof establishes them or guarantees tracking of arbitrary $F^t$.

%======================================================================
\subsection{PL Objective Specialization}
\label{app:pl_specialization}

For smooth nonconvex objectives the result concerns average stationarity; under a Polyak--{\L}ojasiewicz (PL) inequality it also controls the average objective gap.

\begin{corollary}[Average PL objective-gap bound]
\label{cor:app_pl_bound}
Suppose Theorem~\ref{thm:trajrep_stationarity} holds. If $F$ has a global minimizer $w^\ast$ and satisfies $2\mu[F(w)-F(w^\ast)]\leq\|\nabla F(w)\|_2^2$ for all relevant $w$, with $\mu>0$, then
\begin{equation}
\begin{aligned}
\frac1T\sum_{t=0}^{T-1}\mathbb E[F(w^t)-F(w^\ast)]\leq{}&\frac{2[F(w^0)-F_{\inf}]}{\mu\eta aT}\\
&+\frac{2L\eta b_2}{\mu a}\sigma^2\\
&+\frac{2}{\mu a}\left(\frac1{2a}+L\eta\right)\overline{A^2}_T.
\end{aligned}
\label{eq:app_pl_bound}
\end{equation}
\end{corollary}

\begin{proof}
The PL inequality implies $F(w^t)-F(w^\ast)\leq\|\nabla F(w^t)\|_2^2/(2\mu)$. Take expectations, average over effective rounds, and apply Eq.~\eqref{eq:stationarity_main}. Dividing its right-hand side by $2\mu$ gives Eq.~\eqref{eq:app_pl_bound}.
\end{proof}

\noindent This corollary controls the time-averaged gap, not the final iterate; a final-iterate rate needs additional analysis. The benign-descent and harmful-influence assumptions remain necessary.

%======================================================================
\subsection{Finite-Participation Monitoring Recovery}
\label{app:proof_recovery}

\noindent\textit{Proof of Proposition~\ref{prop:finite_time_recovery}.} Let $k_1,\ldots,k_{N_w}$ denote the consecutive valid participations in the proposition, where consecutive means there are no intervening valid participations of this client and any number of global rounds may occur between them. The monitoring threshold may change between these participations. By assumption, $q_i^{k_j}-\delta^{k_j}\geq\varepsilon_q$ for every $j=1,\ldots,N_w$. Since the sigmoid is increasing and $\kappa_{\mathrm{rec}}>0$,
\begin{equation}
\rho_i^{k_j}=\sigma\big(\kappa_{\mathrm{rec}}(q_i^{k_j}-\delta^{k_j})\big)\geq\sigma(\kappa_{\mathrm{rec}}\varepsilon_q)=:\rho_{\min}>\theta_{\mathrm{rec}}.
\label{eq:app_recovery_score_bound}
\end{equation}
At the last of these participations, the window of up to $N_w$ most recent recorded readiness scores contains precisely these $N_w$ scores. Therefore,
\begin{equation}
\frac1{|\mathcal W_i^{k_{N_w}}|}\sum_{k\in\mathcal W_i^{k_{N_w}}}\rho_i^k\geq\rho_{\min}>\theta_{\mathrm{rec}}.
\end{equation}
The recovery condition in Eq.~\eqref{eq:method_recovery} holds, so a client that is still monitored transitions to normal at that participation. If the client already recovered at an earlier participation, the claimed ``no later than'' bound also holds under the previously monitored-state transition rule. This proves the proposition. \hfill$\square$

\noindent This statement concerns the monitoring state only. Recovery has no direct effect on the reliability score, update eligibility, aggregation weight, or global model. In particular, the proposition alone does not imply improved clean accuracy or reduced attack success.
